\par\Needspace{8\baselineskip}\bigskip
\pdfbookmark[3]{Exposition, 2}{pallas-II-article2}
\begin{center}2.\end{center}
\noindent\hspace*{\parindent}Substituting in equation 16 the expression \(\mathrm{d}x^2+\mathrm{d}y^2+\mathrm{d}z^2=\mathrm{d}r^2+\frac{pm\mathrm{d}t^2}{rr}\mathpunct{\raisebox{.45ex}{,}}\) we find \(\frac{\mathrm{d}r^2}{\mathrm{d}t^2}=-\frac{m}{a}+\frac{2m}{r}-\frac{pm}{rr}\mathpunct{\raisebox{.45ex}{;}}\) thus at the points of the orbit where \(r\) is a maximum \SourceLeaf{449}or minimum, we must have \(\frac{m}{a}-\frac{2m}{r}+\frac{pm}{rr}=0\). This gives two values for \(r\), namely \(r=a+\BareRoot{(aa-ap)}\) and \(r=a-\BareRoot{(aa-ap)}\). It is evident from this that \(2a\) is the major axis and \(\frac{\BareRoot{(aa-ap)}}{a}\mathpunct{\raisebox{.45ex}{,}}\) which we shall set \(=e\), is the eccentricity. We thus have
\[
\begin{gathered}
p=a(1-ee)\\[4pt]
\mathrm{d}t^2=\frac{arr\mathrm{d}r^2}{m(aaee-(a-r)^2)}\mathpunct{\cdot}
\end{gathered}
\]
\par To integrate this equation, we set
\[
\begin{array}{rl}17)\quad&a-r=ae\cos E,\end{array}
\]
from which it follows that
\[
\mathrm{d}t^2=\frac{arr\mathrm{d}E^2}{m}\mathpunct{\raisebox{.45ex}{;}}
\]
so, on the permissible assumption that \(E\) always increases, we have
\[
\mathrm{d}t=\BareRoot{\frac{a}{m}}\cdot r\mathrm{d}E=\BareRoot{\frac{a^3}{m}}\cdot(1-e\cos E)\mathrm{d}E.
\]
\par Thus if we set
\[
\begin{array}{rl}
18)\quad&\BareRoot{\dfrac{m}{a^3}}=n\\[6pt]
19)\quad&E-e\sin E-nt=\epsilon,
\end{array}
\]
the equation \(n\mathrm{d}t=(1-e\cos E)\mathrm{d}E\) shows that \(\epsilon\) will be a constant. We see that \(E\) represents what astronomers call the eccentric anomaly, \(nt+\epsilon\) the mean anomaly, \(n\) the mean motion in the unit of time, and \(\epsilon\) the epoch of the mean anomaly at time zero. Further, it is clear that in the time of one revolution \(E\) must increase by \(360^{0}\) or by \(2\pi\), with \(2\pi\) being the circumference of a circle of radius \(1\). Thus if \(T\) denotes this time, we have
\[
nT=2\pi\qquad\text{or}\qquad T=2\pi\BareRoot{\frac{a^3}{m}}\mathpunct{\raisebox{.45ex}{,}}
\]
which amounts to \textsc{Kepler}’s theorem.
\par This also gives us a measure of the masses once we have chosen the units of time and space, namely
\[
m=\frac{4\pi\pi a^3}{TT}\mathpunct{\cdot}
\]
\par Thus, choosing the mean solar day as one unit and the \SourceLeaf{450}value of \(a\) in the Earth’s orbit as the other, and further assuming \(T=365{,}2[563835]\), we find
\begin{center}
\textit{Sum of the masses of the Sun and the Earth}\par\nobreak
\mbox{. . . . \(0{,}000295913[0424]\)}
\end{center}
or, assuming the former to be \(354710\) times as great as the latter
\begin{center}
\textit{Mass of the Sun} . . . . \(0{,}000295912[2082]\),
\end{center}
whose common logarithm is \(6{,}471\;1628\;828-10\).
